\section{Постановка задачи}

\subsection{Контекст исследования}

Как уже было сказано, классические алгоритмы обучения с подкреплением требуют жестко заданной архитектуры нейросети. Чтобы избавиться от этой рутины, целесообразно использовать нейроэволюционный алгоритм NEAT. Имитируя природную эволюцию, он забирает всю лишнюю работу на себя.

\subsection{Цель исследования}

Провести сравнительный анализ классического глубокого Q-обучения и нейроэволюционного подхода.

\subsection{Задачи исследования}

\begin{enumerate}

  \item Реализовать базовую модель обучения с подкреплением для решения задачи моделирования поведения когнитивного агента.

  \item Применить алгоритм NEAT к той же задаче для эволюции архитектуры нейросети агента.

  \item Провести сравнительный анализ моделей по метрикам, определенным ниже.
  
\end{enumerate}

\subsection{Метрики исследования}

\begin{definition}
  \label{def:learning_efficiency}
  Скорость обучения определяет, насколько быстро алгоритм находит решение, удовлетворяющее заданному критерию успеха.
\end{definition}

Из-за различий в процедурах обучения будем использовать две шкалы:

\begin{enumerate}

  \item Количество взаимодействий со средой: Общее число тактов симуляции, потребовавшихся агенту для достижения порогового значения награды $R_{threshold}$. Используем формулу: $N_{steps} = \sum^E_{i=0}T_i$, где $E$ -- количество эпизодов/поколений до достижения успеха, $T_i$ -- длительность i-го эпизода/поколения.

  \item Время обучения: Астрономическое время в минутах, затраченное на обучение до достижения критерия успеха. Используем формулу: $T_{train} = \sum^E_{i=0}t_i$, где $t_i$ -- время i-го эпизода/поколения.
\end{enumerate}

\begin{definition}
  \label{def:quality}
  Качество решения -- эффективность итоговой стратегии агента после завершения обучения.
\end{definition}

\begin{definition}
  \label{def:complexity}
  Сложность модели -- сложность ее обучения с точки зрения вычислений.
\end{definition}

\begin{enumerate}

  \item Количество параметров: Используем формулу: $P = \sum_{l=1}^{L} (n_{l-1} + 1) \times n_l$, где $L$ -- количество слоев, $n_l$ -- число нейронов в l-м слое.

  \item Плотность связей: Используем формулу: $D = \frac{E_{active}}{E_{max}}$, где $E_{active}$ -- число активных связей, $E_{max}$ -- максимально возможное число связей.

\end{enumerate}

\begin{definition}
  \label{def:robustness}
  Устойчивость -- способность агента сохранять ту же производительность при изменении условий среды, на которых он не обучался явно.
\end{definition}

\begin{enumerate}

  \item Устойчивость к шуму наблюдений: Измеряется как процентное снижение средней накопленной награды $\overline{R}$ при добавлении гауссовского шума к входным данным агента. Используем формулу: $Robustness_{noise} = \frac{\overline{R}_{clean} - \overline{R}_{noisy}}{\overline{R}_{clean}} \times 100\%$

  \item Устойчивость к изменениям среды: Измеряется как процентное изменение средней накопленной награды $\overline{R}$ при модификации среды. Используем формулу: $Robustness_{env} = \frac{\overline{R}_{original} - \overline{R}_{modified}}{\overline{R}_{original}} \times 100\%$

\end{enumerate}

\subsection{Математическая формализация}

Среда моделируется как Марковский процесс принятия решений (МППР) $(\mathcal{S}, \mathcal{A}, \mathcal{P}, r, \gamma)$, где $\mathcal{S}$ -- пространство состояний среды, $\mathcal{A}$ -- пространство действий агента, $r : \mathcal{S} \times \mathcal{A} \rightarrow \mathbb{R}$ -- функция награды. Задача обучения с подкреплением формулируется как максимизация ожидаемого кумулятивного вознаграждения:

\begin{equation}
  \label{eq:rl_fitness}
  J(\theta) = \mathbb{E}_{\tau \sim \pi_{\theta}} \left[ \sum^{T}_{t=0} \gamma^t r(s_t, a_t) \right]
\end{equation}

где $\pi_{\theta}$ -- политика агента, параметризованная вектором $\theta$, $\gamma \in [0, 1]$ -- коэффициент дисконтирования, $T$ -- максимальная длина эпизода, $\tau = (s_0, a_0, s_1, a_1, \ldots, s_{T-1}, a_{T-1}, s_T)$ -- траектория взаимодействия агента со средой. Подробнее МППР рассматривается в статьях \cite{Sutton1998} и \cite{mnih2015humanlevel}. В случае текущей постановки МППР мы как раз ссылаемся на \cite{Sutton1998}e

\subsection{Сравниваемые подходы}

\begin{enumerate}

  \item   Классический RL-метод: Оптимизация параметров $\theta$ нейросети фиксированной архитектуры методом градиентного спуска/подъема на основе оценки $\nabla_{\theta} J(\theta)$
  
  \item   Нейроэволюционный подход: Совместная оптимизация топологии сети $\mathcal{T}$ и весов $\theta$ эволюционным алгоритмом, минимизирующим сложность при максимизации $J(\theta, \mathcal{T})$
  
\end{enumerate}

\subsection{Ожидаемые результаты}

Оценка преимуществ и недостатков каждого подхода в соответствии с предложенными метриками, рекомендации по выбору метода в зависимости от характеристик задачи моделирования когнитивных агентов.